% Preamble
\chapter*{Preamble}
\addcontentsline{toc}{chapter}{Preamble}
\label{ch:preamble}

This guide synthesizes two sources of knowledge about using Claude effectively.

\textbf{Official guidance} from Anthropic: documentation, blog posts,
engineering reports, and enterprise case studies---researched and verified
as of February 2026.

\textbf{Practitioner patterns} abstracted from extensive daily use of
Claude Code across 20+ projects spanning data science, machine learning,
software engineering, and technical writing.

Every practice in this guide is labeled:

\begin{description}
  \item[{\color{OfficialPurple}\textbf{[Official]}}] Anthropic's documented recommendation, with source URL.
  \item[{\color{PractitionerTeal}\textbf{[Practitioner]}}] Discovered through extensive daily use.
  \item[{\color{ConvergenceGold}\textbf{[Convergence]}}] Both sources agree independently.
\end{description}

\subsection*{Who This Is For}

\textbf{Primary audience:} Senior ICs and tech leads who already use Claude
and want to level up. This guide assumes familiarity with Claude Code basics.

\textbf{Secondary audience:} Engineering managers evaluating Claude for their teams.
\cref{ch:enterprise} is specifically designed to be shared with decision-makers.

\subsection*{How to Read This Guide}

\begin{description}
  \item[Getting started?] Read \cref{ch:configuration} and \cref{ch:greenfield}.
    Set up CLAUDE.md and your first hooks.
  \item[Already using Claude?] Read \cref{ch:antipatterns} and \cref{ch:maturity}
    to identify your current level and most impactful next investment.
  \item[Evaluating for your team?] Read \cref{ch:enterprise} and \cref{ch:maturity}.
  \item[Optimizing costs?] Read \cref{sec:cost} and \cref{sec:cost-scale}.
\end{description}

\subsection*{Foundational Principles}

Four engineering principles unify every recommendation in this guide:
\emph{never fail silently}, \emph{simplicity over complexity},
\emph{immutability by default}, and \emph{fail fast with diagnostics}.
See \cref{ch:principles} for the full treatment. Keep them in mind
as you read---they explain the ``why'' behind the specific practices.

\subsection*{A Note on Pace}

The AI tooling landscape moves fast. The model capabilities, API features,
and best practices described here are verified as of February 2026.
Some details may have changed by the time you read this.
Check the repository for the latest version, and consult
\cref{app:sources} for primary source URLs.

\vspace{0.5\baselineskip}
\hfill\textit{--- Brandon Behring, February 2026}
